
i.e. qui rend \(\operatorname{Coker}u\) nul. Comme la formation de \(\operatorname{Coker}u\) commute au changement de base, on voit qu'il est représenté par l'ouvert \(S-\lambda(\operatorname{supp}(\operatorname{Coker}u))\) de \(S\). On peut supposer donc que \(u\) est déjà un épimorphisme. Mais utilisant maintenant l'hypothèse que \(G\) est \(S\)-plat, on trouve que la formation de \(\operatorname{Ker}u\) commute au changement de base, et par suite le foncteur envisagé est représenté par l'ouvert \(S-\lambda(\operatorname{supp}(\operatorname{Ker}u))\) de \(S\).



% Expose XIII local macros
\providecommand{\Ass}{\operatorname{Ass}}
\providecommand{\Hom}{\operatorname{Hom}}
\providecommand{\Ext}{\operatorname{Ext}}
\providecommand{\Proj}{\operatorname{Proj}}
\providecommand{\NS}{\operatorname{NS}}
\providecommand{\Tors}{\operatorname{Tors}}
\providecommand{\Picnum}{\operatorname{Pic}^{\mathrm{num}}}
\providecommand{\Piczero}{\operatorname{Pic}^{0}}
\providecommand{\Pictau}{\operatorname{Pic}^{\tau}}
\providecommand{\length}{\operatorname{length}}
\providecommand{\etale}{\text{\'etale}}


